\dykband{EXERCISES IN SWAFL}
{\small Open SWAFL (desktop or online) with the benchmark project. The chapter of the \href{https://github.com/halilunaziru73-creator/Coupled-Water-Fluxes-Pore-Scale-Dynamics-Precision-Irrigation-Mulching-Feedback-Loops/raw/main/Software/v1.1.0/SWAFL_Desktop/docs/manual_v1.1.0/SWAFL_Manual_v1.1.0.pdf}{SWAFL manual} that explains each task is given in brackets; record your values in a table and compare them with what is expected (\emph{Check}).}\par\vspace{3pt}
\begin{exlist}
\item {\color{khblue}[Manual, Chapter 8]} Open \emph{Deterministic phenotyping}; run the well-watered and the drought plants side by side and record the height and leaf area at day 70 and day 140.\par\noindent{\small\itshape Check: the drought plant is smaller and rolls its leaves.}
\item {\color{khblue}[Manual, Chapter 8]} Open \emph{Volumetric phenotyping}; follow the scan step by step and note the number of views and the registration.\par\noindent{\small\itshape Check: the views are joined into one cloud.}
\item {\color{khblue}[Manual, Chapter 8]} Change the soil from the loam to the sand and repeat the phenotyping. Explain the difference with the water held by each soil.\par\noindent{\small\itshape Check: the plant on the sand is stressed earlier.}
\item {\color{khblue}[Manual, Chapter 8]} Change the mulch and the tillage and observe the effect on the plant.\par\noindent{\small\itshape Check: mulch delays the stress.}
\end{exlist}
